\CatchFileBetweenTags{\AlphaInvVal}{calculations/constants.tex}{AlphaInvVal}

\section{Introduction}

Any system that persists by partitioning itself creates a new failure mode due to fragmentation. Subsystems that optimize locally at the expense of global coordination decay. For the partitioned whole to function, information must flow across the boundaries between tiers and this flow must not dissolve the system faster than the partitions protected it.

The cost of boundary crossing is irreducible. Information deforms to match new constraints, works against substrate resistance, and does so under temporal pressure. The product—dissipation density—must be minimized or the system fragments. 

Every boundary crossing has an irreducible cost: the information payload must deform to match the destination's constraints (structural mismatch), overcome the substrate's resistance (impedance), and do so at some rate (flux). The product is dissipation density. Non-optimal transitions accumulate entropy faster than the system can shed it; the system decoheres.

Not all transition geometries are stable. The Selection Principle determines among all possible ways to move information across a partition boundary which ones result in the system persisting. This paper derives the stable morphologies of inter-partition transfer—the geometric attractors that survive the filter.  This applies to cellular organelles, corporate divisions, and lattice sectors alike.

This paper derives the stable geometries of inter-partition information transfer: the morphologies that survive the Selection Principle's filter. We show that only three stable types exist on any finite, discrete substrate, and that a fourth logical possibility—boundary dissolution—triggers catastrophic identity loss.

To validate this architecture, we apply it to the $E_8$ substrate established in \cite{meyer_informational_2026, meyer_ieaction_2026}.

The CKM and PMNS matrices of the quantum vacuum's flavor sector emerge as thermodynamic optima of the three transition types. The structural duality between them—quarks as topological knots constrained by interface friction, neutrinos as lattice phonons exhibiting symmetry transparency—resolves the decades-old puzzle of why these mixing matrices behave so differently under identical unitary mathematics.

\begin{enumerate}
    \item \textbf{The Geometric Determination of Flavor Matrices:} The Geometric Determination of Flavor Matrices: The CKM and PMNS parameters emerge as closed-form geometric aspect ratios, with mixing hierarchies fixed by lattice topology rather than empirical input.
    
    \item \textbf{The Structural Coupling of Gauge and Flavor Sectors:} The Gatto-Sartori-Tonin relation emerges as an impedance-matching condition between the Weak Mixing Angle and the Cabibbo Angle, both partitioning the same finite channel capacity.
\end{enumerate}

\textbf{Scope and Continuity:} All fundamental inputs utilized in these derivations; including the discrete lattice capacity ($\nu=16$), interaction rank ($\sigma=5$), topological boundary ($\chi=2$), spatial manifold ($D=4$), and the vacuum impedance ($\alpha^{-1} \approx \AlphaInvVal$) are taken strictly as rigid priors established in ``Persistence Mechanics: A Universal Architecture of Persistence''\cite{meyer_informational_2026}, ``Persistence Mechanics: Entropic Action''\cite{meyer_ieaction_2026}, ``Persistence Mechanics: Nested Persistence''\cite{meyer_ienested_2026}. \textit{Zero free parameters} are introduced to fit the flavor data.